\begin{titlepage}
\thispagestyle{empty}
\centering
\vspace*{12mm}
{\Huge\bfseries SciScope 科技文献数据要素服务\\与科研智能体系统设计\par}
\vspace{6mm}
{\Large 国赛版架构与完整设计文档（内部）\par}
\vspace{8mm}
{\large \tagpill{数据要素服务} \quad
\tagpill{证据立场 L3} \quad
\tagpill{MCP 证据后端} \quad
\tagpill{本机全量演示}\par}
\vspace{16mm}

\begin{tikzpicture}[
  node distance=8mm,
  every node/.style={font=\small},
  box/.style={draw=black, fill=white, rounded corners=3pt, minimum height=14mm, align=center},
  arrow/.style={-{Latex[length=2.2mm]}, thick, color=black}
]
  \node[box, minimum width=34mm] (data) {公开科技文献\\六源采集 / 数据要素治理};
  \node[box, minimum width=37mm, right=10mm of data, very thick] (analysis) {可计算数据要素\\分析资产 / RAG 语料 / 图谱};
  \node[box, minimum width=37mm, right=10mm of analysis, very thick] (agent) {证据层智能体\\检索 / 问答 / 趋势 / 推荐\\stance 判定 / 争议地图};
  \node[box, minimum width=48mm, below=15mm of agent, double] (mcp) {协议出口\\API / SSE / MCP（tools + resources）\\被其他科研 agent 调用};
  \node[box, minimum width=34mm, left=10mm of mcp] (tui) {展示端\\Go TUI 消费 SSE};
  \draw[arrow] (data) -- (analysis);
  \draw[arrow] (analysis) -- (agent);
  \draw[arrow] (agent) -- (mcp);
  \draw[arrow] (agent) -- (tui);
\end{tikzpicture}

\vfill
\begin{designbox}{文档定位}
本文档是团队内部用于统一国赛架构、证据层（stance L3）、协议边界、数据与演示路线的
设计文档，是 `docs/project/国赛目标说明书.md`（冻结目标）在设计层的展开。
对外交付仍落在两条主线上：可复现的数据分析报告（要素服务化表述），以及可被调用、
可演示的科研智能体系统。核心判断：**证据可追溯是魂，协议可调用是门，界面可替换是皮**。
\end{designbox}
\vspace{6mm}
{\small 版本: 国赛版（2026-08-04） \quad 上一版: 2026-06-19}
\end{titlepage}
\clearpage
